\chapter{豌豆为什么让遗传变得可计算}

\begin{kaichang}
打开冰箱冷冻层，你家也许也有一袋冷冻豌豆。拿一粒放在手心：小小的，圆圆的，绿绿的。别忘了，它不是"食品厂的零件"，而是一个"孩子"——豌豆植株受精之后结出的种子。只要把它埋进湿润的土里，它就能长成一整棵豌豆苗，再开花，再结出一荚新的豌豆。

现在请你猜一个问题。一株结圆粒豌豆的植株，和一株结皱粒豌豆的植株杂交——植物的"婚配"——它们的孩子会结什么样的种子？圆的？皱的？还是"半圆半皱"，像两种颜料兑在一起调出的折中颜色？

把你的猜测写在纸上，这一章结束时我们回来对答案。一百六十多年前，一位名叫孟德尔的修道士，就是在修道院的菜园里种了两万多株豌豆，把这个问题从"看相"变成了"算术"。为什么偏偏是不起眼的豌豆，让遗传第一次变得可以计算？这正是本章的标题，也是我们这一路要拆开的谜。
\end{kaichang}

\section{先看清现象：孩子像父母，但又不是"兑"出来的}

先来给本章的主角起个人话名字。\textbf{性状}，就是生物身上那些"看得出来、数得出来"的特征：豌豆种子是圆是皱、子叶是黄是绿、植株是高是矮；放在人身上，就是能不能把舌头卷成筒、耳垂是贴着还是垂着、大拇指最后一节能不能向后弯、笑起来有没有酒窝。你现在就可以对照镜子数一数自己——不用任何仪器，这些都是肉眼可见的。

有些性状是"非此即彼"的：舌头要么卷得起来，要么卷不起来，很少有人"卷一半"。有些性状却是连续的：身高可以从一米四排到两米，肤色有无数种深浅。这个区别看似不起眼，却几乎决定了遗传学能不能被发明——记住它，等下孟德尔正是靠这一点"作弊"的。

在孟德尔的时代，人们对遗传有一个听起来天经地义的默认理论：\textbf{融合遗传}。父母各出一半"特质"，像一滴红墨水滴进一滴蓝墨水，孩子就是兑出来的紫色。它和日常经验似乎很合拍：孩子的身高常常介于父母之间，肤色也常常像父母的"中间色"。

可这个理论藏着一个致命的预言，很少有人认真推演过：如果遗传真是混合，那么每繁殖一代，差异就被稀释一半。高个子配矮个子出中等个，中等个配中等个还是中等个——用不了多少代，所有人都会挤在平均值附近，世界的变异会被"搅匀"直至消失。可现实呢？变异从来没有消失：平均身高的父母照样生出特别高或特别矮的孩子，两个卷舌的人能生出不会卷舌的孩子，爷爷奶奶的特征会跳过一代忽然在孙辈身上"复活"。

这说明"特质"不是墨水，混合不掉。但光凭这种生活观察，你永远说服不了对手——"差不多""好像""我觉得"吵不出结果。要一锤定音，需要把遗传变成可以计数的实验：找一种生物，选几对性状，让亲本交配，然后\textbf{一株一株地数}。

\section{孟德尔的菜园：把"看相"变成"数数"}

格雷戈尔·孟德尔，1822 年出生在一个贫苦农家，后来进了奥地利布隆（今捷克布尔诺）的一座修道院。他不是"灵机一动"的天才：他在维也纳大学正经学过数学和物理，物理课上学的那套"设计对照、精确测量、用数字说话"的手艺，全被他搬进了菜园。从 1856 年起，他花了大约八年，做了一件前无古人的事。这件事的每一步都精心设计过，值得逐一拆开看。

\textbf{第一步，选材料：为什么是豌豆？}豌豆简直是为此而生的实验生物。它自花授粉——一朵花的雄蕊和雌蕊被花瓣包着，天然"自己跟自己结婚"，所以世世代代容易保持纯正，可以养出性状稳定不漂移的"纯系"；它的花又足够大，花药可以人工剪除，再拿毛笔蘸上另一株的花粉涂上去，交配对象完全由实验者指定；它生长快，一年就能收获下一代；每株能结几十上百粒种子，后代数量管够；最妙的是，种子本身就是下一代个体——数豆粒，就是在数"孩子"。

\textbf{第二步，选性状：专挑"非此即彼"的。}孟德尔没有贪多，只挑了七对界限分明的性状：种子圆或皱、子叶黄或绿、花紫或白、豆荚饱满或缢缩、豆荚绿或黄、花腋生或顶生、茎高或矮。连续变化的性状他一律不碰。这是他最聪明的决定之一：只有"非此即彼"，计数才没有含糊，比例才会从数据里自己站出来。

\textbf{第三步，控制交配。}在豌豆花还没成熟时剪掉雄蕊（去雄），套上纸袋防止"第三者"的花粉闯入，再人工授上指定亲本的花粉。每一株孩子的父母是谁，全部记录在案，没有一笔糊涂账。

\textbf{第四步，逐代、大量地计数。}纯系亲本杂交，得到的第一代孩子（记作 $F_1$）全部留下种子；$F_1$ 再自交，得到第二代（$F_2$），继续数。八年下来，他检查的豌豆植株超过两万八千株。

结果的第一幕就让人意外。圆粒纯系和皱粒纯系杂交，$F_1$ 的种子\textbf{全是圆粒}，皱粒性状像被整个吞掉了——融合遗传在这里似乎赢了一局：一种性状"覆盖"了另一种。但第二幕剧情反转：让 $F_1$ 自交，皱粒在 $F_2$ 里原封不动地回来了，不多不少，大约占四分之一。孟德尔数出来的是 5474 粒圆、1850 粒皱，比值约 $2.96:1$。其余六对性状如出一辙：子叶黄绿 $6022:2001$，花色紫白 $705:224$，茎的高矮 $787:277$——每一对的 $F_2$ 都紧贴 $3:1$。

这两个现象合在一起，是通往答案的钥匙：皱粒性状没有在 $F_1$ 里被"稀释"，否则它不可能在 $F_2$ 里完整地、以精确的比例复活。它在 $F_1$ 身体里只是被\textbf{遮盖}了，原样保存着，一代之后又跳了出来。遗传的载体不是墨水，而是一颗一颗、传来传去不会被磨掉的"颗粒"。

\section{粒子层：给看不见的"因子"安排座位}

看见比例只是第一步，解释比例才是科学。孟德尔为这些"颗粒"设计了一幅极简的内部图景，只用了三条假设：

\begin{itemize}[itemsep=2pt]
\item 每株豌豆的每个性状，由\textbf{一对}遗传因子掌管，两个因子一个来自父本、一个来自母本。
\item 形成配子（花粉里的精细胞、胚珠里的卵细胞）时，这一对因子\textbf{彼此分开}，每个配子只带走一个。
\item 受精时，两个配子各交出一个因子，\textbf{随机}配成新的一对。
\end{itemize}

这几句话今天写在课本里平平无奇，在当时却是石破天惊：从来没有人敢假设生物体内有这样成对的、会分开、会重新抽签组合的离散单元。后来的人给这些单元换上了今天的名字。决定某个性状的那个"位置"叫\textbf{基因}（gene），一个基因的不同"版本"叫\textbf{等位基因}（allele）——圆粒版和皱粒版，是同一个基因的两个等位基因。一个个体的两个等位基因组合，叫\textbf{基因型}（genotype）；这对组合最终呈现出来的样子，叫\textbf{表型}（phenotype）。两个版本相同（圆圆或皱皱）叫\textbf{纯合}，不同（圆皱各一）叫\textbf{杂合}。

还有一个必须马上说清的概念。\textbf{显性}与\textbf{隐性}：杂合个体身上只显现其中一个版本的效果，显现出来的那个叫显性（圆粒），被遮盖的叫隐性（皱粒）。请注意这句话的措辞——显性隐性描述的是"杂合子的表型长什么样"，仅此而已。它不是两个因子在体内打了一架、赢的说了算；因子不会打架。显性往往是分子层面的互作结果：比如一个版本能造出有功能的产物、另一个几乎造不出来，而一份产物常常就够用了，杂合子便呈现显性表型。基因怎样通过蛋白质这类化学物质影响性状，中间隔着一整套生命过程，第 73 章会专门讲。这里你只需记住：显隐性是"产物的供求关系"，不是"因子的强弱排名"。

\section{规则层：为什么是 $3:1$，而不是别的数}

有了粒子层的图景，$3:1$ 就不再是巧合，而是一道算术题。设圆粒等位基因为 $R$，皱粒为 $r$。$F_1$ 全是杂合子 $Rr$。它形成配子时，一半配子带 $R$，一半带 $r$——就像抛一枚均匀的硬币。$F_1$ 自交，相当于两枚硬币同时抛：

\begin{qiaoliang}
每次受精，父本配子带 $R$ 或 $r$ 的概率各为 $\frac{1}{2}$，母本配子亦然。两次抛硬币互相独立，第 8 章讲过：独立事件同时发生的概率等于各自概率相乘。于是四种组合等概率：
\[
RR:\ \tfrac{1}{4},\qquad Rr:\ \tfrac{1}{4},\qquad rR:\ \tfrac{1}{4},\qquad rr:\ \tfrac{1}{4}.
\]
$Rr$ 和 $rR$ 是同一基因型，所以基因型之比为 $RR:Rr:rr = 1:2:1$。又因为 $RR$ 和 $Rr$ 的表型都是圆粒，表型之比为
\[
\text{圆粒}:\text{皱粒} = \tfrac{3}{4}:\tfrac{1}{4} = 3:1.
\]
这就是\textbf{分离定律}：成对的等位基因在形成配子时彼此分离，各自进入不同的配子。$3:1$ 不是豌豆的癖好，而是"$1:2:1$ 的基因型比"加上"显性遮盖"之后的必然外观。
\end{qiaoliang}

这套逻辑马上解释了更多数字。孟德尔接着把两对性状放在一起跟踪：圆粒黄子叶（$RRYY$）配皱粒绿子叶（$rryy$），$F_1$ 全是圆黄（$RrYy$），$F_1$ 自交后，$F_2$ 出现四种表型：圆黄、圆绿、皱黄、皱绿，比例约 $9:3:3:1$。这个数也藏不住来历：圆皱之比 $3:1$，黄绿之比也是 $3:1$，而两对性状的分离互不相干，所以组合概率相乘——$(3+1)\times(3+1)$ 摊开正是 $9:3:3:1$。这就是\textbf{独立分配定律}（自由组合定律）：不同对的等位基因各自独立地分离、自由地组合。先记住它字面的样子，第 64 章我们会看到它何时成立、何时失效——基因们并不总是这么"互不相干"。

你可能想问：凭什么数出来的是 $2.96:1$ 而不是精确的 $3:1$？答案是概率的本性。抛四次硬币，"两正两反"很常见，"三正一反"也不稀奇；只生四个孩子，性别恰好两男两女的概率其实只有 $\frac{3}{8}$。样本越小，偏离越大；样本越大，比例越被大数定律按向理论值。这正是孟德尔要数几千株的深层原因——两万八千株豌豆，平均每年约三千五百株，摊到一百多天的生长季里，每天要授粉、套袋、挂牌、观察几十株。$3:1$ 这个优美的数字，是整整八个春天一株一株数出来的。

概率还能反过来用，解决一类"看表型猜基因型"的问题：从 $F_2$ 的圆粒种子里随机挑一粒，它体内藏着皱粒等位基因的概率是多少？圆粒里 $RR$ 和 $Rr$ 之比是 $1:2$，所以答案是 $\frac{2}{3}$。这类"已知结果、反推原因"的算法叫条件概率，看起来只是算术游戏，实际上正是今天遗传咨询的基本功——本章的挑战层还会用到它。

\section{符号层：字母终于登场}

到此为止我们故意一个遗传符号都没用，现在它们可以出场了，因为你已经知道每一笔账的含义：

\begin{itemize}[itemsep=2pt]
\item 显性等位基因用大写字母（$R$），隐性用同一字母的小写（$r$）。大小写纯粹是人类的记账约定，不代表字母本身有大小强弱。
\item 亲本记作 $P$，子一代 $F_1$，子二代 $F_2$；"$\times$"表示杂交，"$F_1$ 自交"表示一株植株自己授粉给自己。
\item 配子类型写在箭头两端：$Rr \rightarrow$ 配子 $R$、$r$。
\end{itemize}

最顺手的记账工具是\textbf{棋盘格}（庞尼特方格）：把父本、母本的配子分别写在格子的顶部和左侧，交叉格里就是受精卵的基因型，每格概率相等。$Rr \times Rr$ 的棋盘格如下：

\begin{table}[htbp]
\centering
\caption{分离定律的棋盘格：每一格占 $\frac{1}{4}$}
\begin{tabular}{c|cc}
\toprule
 & 配子 $R$ & 配子 $r$ \\
\midrule
配子 $R$ & $RR$（圆） & $Rr$（圆） \\
配子 $r$ & $Rr$（圆） & $rr$（皱） \\
\bottomrule
\end{tabular}
\end{table}

一眼就能读出：三格圆、一格皱。画棋盘格的能力比背"$3:1$"重要得多——背下来的比例会忘，画出来的格子永远能重新算。

同样，概率树是另一种记账法：从"$F_1$ 产生配子"分出两支（各 $\frac{1}{2}$），每支再分两支，沿着树枝把概率一路乘下去，末端就是各种合子的概率。棋盘格和概率树是同一枚硬币的两面：一个用面积，一个用路径，记的都是同一本账。遇到两对性状，棋盘格要画成 $4\times4$ 的十六格（配子为 $RY$、$Ry$、$rY$、$ry$），格子变多，逻辑不变——你不想画格子的时候，就用乘法：$(3:1)\times(3:1)=9:3:3:1$。

最后提醒一句：$R$、$r$ 只是占位符号，它们不代表等位基因"长什么样"。等位基因的物理本体是什么、住在细胞的什么地方，孟德尔一概不知道，也不影响他的算术成立——这是科学史上罕见的"先把规律算对、再找到实体"的案例。

\section{证据层：怎样让怀疑者闭嘴}

\begin{zhengju}
一套漂亮的假设不算科学，能经受"换着花样来推翻"才算。孟德尔自己也清楚这一点，他为"成对因子分离重组"的假说安排了最严厉的检验——\textbf{测交}。

怀疑者可以提出一个替代解释：也许 $F_1$ 的皱粒性状不是"藏起来"，而是真的被消灭、被改造了，$F_2$ 里那四分之一皱粒是重新长出来的新东西。怎么区分"藏起来"和"被消灭"？孟德尔让 $F_1$（$Rr$）和皱粒纯系（$rr$）杂交。如果 $F_1$ 体内真的完整保存着一份 $r$，它就会产生等量 $R$、$r$ 两种配子，与 $rr$ 的配子结合后，后代应当是圆粒、皱粒各一半，即 $1:1$；如果 $r$ 已被消灭或稀释，这个比例就不该出现。实验结果：后代圆皱之比紧贴 $1:1$。预言命中——$F_1$ 体内确实藏着一份原封未动的皱粒因子。这个实验的漂亮之处在于：它不是再"看"一遍现象，而是从假说里推出一个\textbf{全新的、此前没人见过的预言}，再让实验来投票。

然而接下来的剧情并不励志。1865 年孟德尔在地方自然学会宣读论文，1866 年正式发表，寄给了当时不少知名学者——然后，几乎没有回音。学界的主流是融合遗传，一个用数学处理生物学的修道士，写的又是"看不见的假想因子"，多数人要么读不懂，要么不以为然。据说著名植物学家内格里收到他的论文后，回信建议他"再研究一种植物"——山柳菊。孟德尔照做了，结果一败涂地：山柳菊行无融合生殖，后代根本不按分离定律走，多年心血几乎动摇了他的信心。这桩公案后来被理解为一段令人惋惜的插曲：不是孟德尔错了，是他被建议研究的材料恰好不适用他的定律。

转机出现在 1900 年——孟德尔去世十六年后。三位植物学家，荷兰的德弗里斯、德国的科伦斯、奥地利的切尔马克，在各自独立做杂交实验、各自得出类似结论后，检索文献时撞见了那篇沉睡三十五年的论文，不得不集体承认：优先权属于孟德尔。同年稍后，贝特森在英国大力宣讲，"遗传学"（genetics）这个词随之诞生。一项发现被遗忘三十五年还能原地复活，靠的不是运气，而是论文里那些可以原样重复的数字。

这个故事还有一个更耐人寻味的尾声。1936 年，统计学家费希尔重新分析孟德尔的原始数据，发现一个"甜蜜的烦恼"：数据\textbf{太好了}，好到与理论值的贴合程度超过统计学允许的自然波动——按概率，真实实验不该这么齐整。学界至今对此有不同猜测：也许助手凭直觉剔除了"长相可疑"的植株，也许数到接近预期就收工。但请注意这场争论的性质：后人无数次重复豌豆及其他生物的杂交，$3:1$ 始终成立，定律本身岿然不动；被质疑的是记录过程的细节。科学就这样工作——它允许对英雄的数据吹毛求疵，因为规律不靠英雄的名望担保，靠可重复性担保。
\end{zhengju}

\section{边界层：什么时候 $3:1$ 会失灵}

孟德尔定律是强大的，但不是宇宙的通用语法。它有明确的适用条件，每一条都值得记成一句"如果……就……"：

\begin{itemize}[itemsep=2pt]
\item \textbf{二倍体、有性生殖、细胞核基因}才谈得上"成对分离"。细菌只有一个环状 DNA，孟德尔比例无从谈起；线粒体和叶绿体的基因几乎只随母本传递，也完全不按此规则出牌（第 66 章附近会再见它们）。
\item \textbf{完全显性}时 $F_2$ 表型才是 $3:1$。显性关系一旦不"完全"，表型比例就会变脸——但注意，基因型的 $1:2:1$ 从不变脸，变的只是表型的外观（下一节细说）。
\item \textbf{受精随机、各基因型存活率相当}。若某种基因型的胚胎难以存活，比例立刻偏斜——某些隐性等位基因纯合致死时，$F_2$ 会出现 $2:1$ 这种"怪比例"，它不是定律的反例，而是定律加上"存活筛选"后的推论。
\item \textbf{样本足够大}。四个孩子里没有出现隐性性状，什么也不能说明；概率只对大样本承诺比例。
\item \textbf{基因之间真正独立}。独立分配定律要求两对基因位于不同染色体上、或在同一条染色体上相距足够远。若它们紧挨在一起，就会"手拉手"一起传递，$9:3:3:1$ 随之走样——这叫连锁，第 64 章讲染色体时会细讲。有趣的是，豌豆恰好只有七对染色体，而孟德尔选的七对性状大多落在不同染色体或相距很远的位置上，他的"好运气"里藏着选择的智慧。
\item \textbf{性状基本由一对基因主控}。身高、肤色、产量这类性状由几十上百个基因共同决定，叠加环境的推拉，根本不会出现 $3:1$——它们是第 73 章的主角。
\end{itemize}

所以请把孟德尔定律当作一张精密的地图：在"单基因、核基因、大样本"的疆域里毫厘不差；走出疆域，它不会错得离谱，而是会以一种你能解释的方式变形。真正高明的使用者不是背下地图，而是知道地图的边界画在哪里。

\section{当表型不再"非此即彼"}

边界层说了"$3:1$ 会变脸"，现在把脸翻开看。这些情形不是对孟德尔的推翻，而是对孟德尔的\textbf{补充}——因为底层基因型的分离比纹丝不动。

\textbf{不完全显性}。紫茉莉（以及金鱼草）的红花纯系与白花纯系杂交，$F_1$ 既不红也不白，而是粉红；$F_1$ 自交，$F_2$ 开出红、粉、白三种花，比例恰好 $1:2:1$。发生了什么？杂合子只造得出一份红色素原料，颜色便"打了对折"——两个等位基因谁也压不住谁，表型把基因型直接"显示"了出来。这种情形其实是观察者的礼物：表型比第一次等于基因型比，孟德尔的 $1:2:1$ 第一次能被肉眼直接看见。

\textbf{共显性}。人的 ABO 血型里，$I^A$ 等位基因让红细胞表面挂 A 型糖链，$I^B$ 挂 B 型糖链；$I^A I^B$ 的杂合子不是"A 和 B 的折中"，而是两种糖链\textbf{同时}挂在细胞表面——血型为 AB 型。两个版本都完整表达，谁也不遮盖谁，这叫共显性。

\textbf{多等位基因}。同样在这个血型系统里，人群中流转着 $I^A$、$I^B$、$i$ 三个版本（$i$ 对前两者都是隐性）。但请盯紧这句话：多等位基因是\textbf{群体}层面的丰盛，不是个体层面的——每个人身体里那个基因座上，仍然只有两个等位基因，一个来自父亲，一个来自母亲。O 型是 $ii$，AB 型是 $I^A I^B$，A 型可能是 $I^A I^A$ 或 $I^A i$。六个基因型、四种表型，全部能用同一个棋盘格算清楚。

\begin{wuqu}
"显性等位基因就是更强、更常见、更有利的版本。"——这个说法流传极广，而且三个词全错。

\textbf{更常见？}O 型血（$ii$）是隐性，却是全世界很多地区最常见的血型；卷舌能力是显性，可不会卷舌的人照样一抓一大把。一个等位基因在人群中常见与否，取决于演化的历史——迁徙、偶然、选择——跟它是显性还是隐性毫无关系。等位基因的频率是群体层面的事，第 85 章前后讲进化时我们还会细算。

\textbf{更有利？}亨廷顿病是一种晚发的神经退行性疾病，致病等位基因恰恰是显性的——只要带一份，人到中年就可能发病；而白化病、苯丙酮尿症这些通常由隐性等位基因引起。"显性"两个字里没有任何关于好坏的承诺。

\textbf{更强？}显隐只描述杂合子的表型长什么样，是分子层面"产物够不够用"的供求关系，不是两个等位基因掰手腕。何况同一个等位基因，在这个组织里是显性，换个环境、换个性状，表现就可能不同。

一句话记住：显性 $\neq$ 强大，显性 $\neq$ 常见，显性 $\neq$ 优秀。下次在广告或传言里听到"显性基因更优秀"，你就有反例可用了。
\end{wuqu}

\begin{shiyan}
\textbf{手心里的分离定律}（约 30 分钟）。

材料：两枚一元硬币（或从两副扑克牌里各抽一红一黑两张牌）、纸和笔。

步骤：把两枚硬币分别命名为"父本"和"母本"。约定正面代表显性等位基因 $R$，反面代表隐性 $r$——这样每枚硬币就是一个杂合亲本 $Rr$ 的"配子发生器"。两枚同时抛出，记录组合：两正记 $RR$，一正一反记 $Rr$（不用区分正反顺序），两反记 $rr$；其中 $RR$ 和 $Rr$ 记"显性表型"，$rr$ 记"隐性表型"。重复抛 40 次，统计基因型比例和表型比例。然后和同桌合并数据，凑到 80 次、160 次，再各算一次比例。

观察点：只抛 12 次时，你很可能得到离 $1:2:1$ 相当离谱的结果；次数越多，比例越被"按"向理论值。这不是你的手法变准了，是大数定律在你手心里发生——孟德尔数两万八千株豌豆，对付的正是同一种运气。

安全提示：硬币、豆子等小物件请远离幼儿和宠物，防止误吞；如果改用豆子做实验，摸过多次、落地过的豆子就不要再入口了。最后郑重提醒：这个游戏模拟的是抽象的概率机制，请不要拿它去给家人的任何性状"算命"或评判任何人。
\end{shiyan}

\section{再看那袋冷冻豌豆}

回到开场那粒豌豆，现在你可以给自己讲完整的故事了。

为什么圆粒纯系配皱粒纯系，孩子不是"半圆半皱"？因为性状不是颜料，是颗粒：$F_1$ 的基因型是 $Rr$，圆粒等位基因的产物够用，皱粒版本被遮盖，于是整荚 $F_1$ 种子全是圆的。为什么皱粒在 $F_2$ 又"复活"了？因为它从未消失——它在 $F_1$ 体内原样保存，分离、重组之后，四分之一的合子拿到 $rr$，皱粒堂堂正正地回来了。你在纸片上写的猜测如果对上了，恭喜你独立走到了 1865 年；如果没对上，更好，你现在知道了为什么。

还有一个彩蛋，把这一章悄悄接回全书的化学主线。皱粒豌豆为什么皱？现代分子生物学找到了答案：皱粒等位基因里插进了一段多余的序列，一种催化淀粉分支的酶因此失灵，种子里糖多、淀粉少，成熟时就因失水而皱缩。你看，等位基因的"显与隐"，追到底是一种酶的工作状态；遗传信息的差异，最终是\textbf{化学组成的差异}。这条"信息怎样变成物质"的暗线，会在第 66 到 69 章彻底展开。

顺带一提，市售的甜脆豌豆常常带有皱粒血统——糖多、淀粉少正是甜的来源。你炒饭里的那口甜，是孟德尔 $F_2$ 里那四分之一的遗产。

\section{继续深挖：从菜园到育种田和医院}

孟德尔的算术早已走出菜园。

在育种田里，它是预测工具。育种家让两个各有优点的亲本杂交时，心里清楚：显性性状在 $F_1$ 就能看到，被遮盖的隐性优点要等到 $F_2$ 才会以约四分之一的频率露面——所以要种够大的群体，才捞得到想要的组合；杂交水稻、抗病小麦的选育流程，账本上都写着 $3:1$ 和 $9:3:3:1$ 的影子。

在医院里，它是风险算术。苯丙酮尿症、白化病等通常由隐性等位基因引起的疾病，携带者（$Aa$）本人完全健康，而当父母双方恰好都是携带者时，每个孩子有约 $\frac{1}{4}$ 的概率患病、$\frac{1}{2}$ 的概率成为和自己一样的携带者。遗传咨询师的工作，就是本章那套条件概率的严肃版本。不过请守住本书一贯的边界：这些都是\textbf{群体层面}的概率，对任何具体的家庭，遗传规律给不出"会不会"的判决书，只给得出"多大可能"的参考；真有疑问，请找正规医院的专业人士，本书不提供任何个体诊疗建议。

最后，把本章最大的悬案留给你。孟德尔把因子算得明明白白，却至死不知道它们\textbf{住在哪里、由什么构成}——他笔下的 $R$ 和 $r$ 只是账簿上的符号，没有实体。他去世前后，显微镜技术突飞猛进，细胞学家在细胞核里看到了一种会在分裂时跳舞的丝状物：染色体。它们成对存在，分裂时彼此分开，配子形成时恰好减半，受精时恰好恢复成对——舞步和孟德尔因子的分离、组合惊人地平行。这是巧合，还是因子就骑在染色体上？如果是，那同一条染色体上的两个基因还怎么"独立分配"？下一章，我们把镜头推进细胞核，看染色体怎样把孟德尔的算术搬进细胞——顺便看它在哪儿摔了一跤。

\begin{tiaozhan}
\textbf{反推的艺术：从表型算回基因型。}一种由隐性等位基因 $a$ 引起的罕见疾病，只有 $aa$ 会发病。一对健康的夫妇，各自都有一个患病的兄弟姐妹（且双方的祖父母辈信息无从得知）。问：他们的第一个孩子患病的概率是多少？

拆解如下。某人的兄弟姐妹患病（$aa$），说明其父母必然都是携带者 $Aa \times Aa$。而此人本人健康，就排除了 $aa$——在"健康"这个条件下，他是 $AA$ 的概率为 $\frac{1}{3}$，是 $Aa$ 的概率为 $\frac{2}{3}$（回忆正文圆粒种子的 $\frac{2}{3}$，同一个结构）。夫妇二人互为独立事件，都是携带者的概率为 $\frac{2}{3}\times\frac{2}{3}=\frac{4}{9}$；两个携带者的孩子患病概率为 $\frac{1}{4}$。相乘得
\[
\frac{2}{3}\times\frac{2}{3}\times\frac{1}{4}=\frac{1}{9}\approx 11\%.
\]
注意每一步都只是一次条件概率，难的不是算术，而是想清楚"哪一步该剔除哪个选项"。这正是遗传咨询的核心算法。跳过本节不影响主线。
\end{tiaozhan}

\section*{练一练}

\begin{sikaoti}
\item 杂合高茎豌豆（$Tt$）与矮茎豌豆（$tt$）测交。画出棋盘格，预言后代的基因型比例与表型比例，并用一句话说明：这个实验为什么能"看见"$F_1$ 体内藏着什么。
\item $Rr$ 自交，只收获了 4 粒种子。恰好 3 粒圆、1 粒皱的概率是多少？（提示：画概率树或列出四种位置，每粒圆粒概率 $\frac{3}{4}$。）算出数字后回答：它离 $100\%$ 有多远？这说明为什么孟德尔不能只种四株？
\item 紫茉莉的红花（$RR$）与白花（$rr$）杂交得粉花 $F_1$，$F_1$ 自交。画出棋盘格，写出 $F_2$ 的基因型比与表型比，并解释：为什么在这里表型比第一次"如实显示"了基因型比？
\item 画一张信息流图：从"亲本体细胞中的成对等位基因"出发，经过"配子"，到"合子"，再到"表型"。在图上标出：哪一步发生了分离？哪一步发生了随机组合？哪一步决定了显性隐性"显形"？
\item 父亲 A 型血、母亲 B 型血，却生出了 O 型血的孩子。写出父母双方可能的基因型，画棋盘格说明 O 型孩子从哪里来；再回答：这对夫妇下一个孩子是 AB 型的概率是多少？
\item 你的同桌说："卷舌是显性性状，所以再过几十代，全世界就人人都会卷舌了。"请写出他的推理里偷换了哪两个概念，并举出本章的两个反例来反驳。
\end{sikaoti}
